% Generated by reproducibility.historical_evaluation_assets; do not edit.
% Accepted source run(s): smartdca-historical-study-v1-5b10a2aba05f84eacfef87b421a580cf7c0dc30d2844c51be6241bc682e39221
% Accepted source artifact digest: cb79fcb63ecb2c9a0894a8f24255dfd4934eaa0dcb54fb4f0894d1b71b941718

\begin{table}[htbp]
  \centering
  \scriptsize
  \setlength{\tabcolsep}{3.2pt}
  \renewcommand{\arraystretch}{1.12}
  \resizebox{\textwidth}{!}{%
  \begin{tabular}{llrrrrl}
    \hline
    Tier & Comparison & Cells & Negative & Zero & Positive & Median range / role \\
    \hline
    H1 & Corrected--DCA & 18 & 18 & 0 & 0 & -4.593\% to -0.335\%; registered complete system \\
    H2 & Corrected--neutral & 18 & 17 & 0 & 1 & -0.545\% to +0.052\%; registered signal only \\
    S1 & Neutral--DCA & 18 & 18 & 0 & 0 & -4.365\% to -0.340\%; descriptive architecture only \\
    \hline
  \end{tabular}%
  }
  \caption{Comparison-tier separation for the eighteen non-unit primary frictionless cells. Corrected--DCA uses DCA wealth as denominator; corrected--neutral uses neutral guarded wealth; neutral--DCA uses DCA wealth. H1 and H2 share the registered Holm family, whereas S1 has no significance decision and is not a causal decomposition. Complete-system significance is never transferred to the signal-only or architecture-only tier.}
  \label{tab:historical-comparison-tiers}
\end{table}

\begin{table}[htbp]
  \centering
  \scriptsize
  \setlength{\tabcolsep}{3.2pt}
  \renewcommand{\arraystretch}{1.12}
  \resizebox{\textwidth}{!}{%
  \begin{tabular}{lrrrr}
    \hline
    Guarded policy & Cells & Mean cash drag & Mean asset exposure & Mean floor activation \\
    \hline
    Corrected & 18 & 2.126\%--11.482\% & 91.698\%--98.909\% & 3.990\%--100.000\% \\
    Neutral & 18 & 1.667\%--8.329\% & 92.041\%--99.563\% & 3.900\%--100.000\% \\
    \hline
  \end{tabular}%
  }
  \caption{Policy mechanisms across the eighteen non-unit primary frictionless cells. Cash drag is terminal policy cash divided by contributed deposits, asset exposure is terminal asset value divided by that policy's cash-inclusive terminal wealth, and floor activation is the mean share of purchase dates with an active clipped guardrail floor. The corrected row comes from H1 left-policy ledgers and the neutral row from S1 left-policy ledgers; these mechanism ranges are descriptive and do not determine a performance ordering.}
  \label{tab:historical-policy-mechanisms}
\end{table}

\begin{table}[htbp]
  \centering
  \scriptsize
  \setlength{\tabcolsep}{3.4pt}
  \renewcommand{\arraystretch}{1.12}
  \resizebox{\textwidth}{!}{%
  \begin{tabular}{lrrrr}
    \hline
    Comparison population & Cells & Lowest 5\% gap & Largest shortfall & Mean cash / valued-unit signs \\
    \hline
    Corrected--DCA, primary frictionless & 18 & -15.292\% & 18.961\% & 18 positive / 18 negative \\
    \hline
  \end{tabular}%
  }
  \caption{Downside and terminal-inventory attribution for the corrected guarded policy against DCA. The 5\% value is the lowest linearly interpolated cellwise quantile, and largest shortfall is the maximum observed episode loss magnitude. Mean cash contribution is $C^c-C^D$ and the valued-unit contribution is $P(Q^c-Q^D)$; their sum is the mean dollar wealth gap in every cell. The signs describe ledger-conditioned accounting, not a causal explanation.}
  \label{tab:historical-risk-attribution}
\end{table}
\FloatBarrier
